\documentclass[11pt]{article}
\usepackage[margin=1.1in]{geometry}
\usepackage{amsmath,amssymb,amsthm}
\usepackage{booktabs}
\usepackage{listings}
\usepackage[colorlinks=true,linkcolor=blue,urlcolor=blue]{hyperref}
\usepackage{microtype}

\lstset{basicstyle=\ttfamily\small,breaklines=true,frame=none,
        columns=fullflexible,xleftmargin=1em,aboveskip=6pt,belowskip=6pt}

\newtheorem{definition}{Definition}[section]

\title{Reinforcement learning environments for program verification:\\
what a proof kernel does and does not give you}
\author{Youssef \textsc{Miled}}
\date{August 2026}

\begin{document}
\maketitle

\begin{abstract}
\noindent
Training language models to verify programs is an attractive proposition because
a proof assistant supplies a reward that cannot be fooled: a proof either
type-checks against the kernel or it does not. This is unlike the rewards
normally used for code generation, where a test suite passing is only weakly
related to a program being correct. It is often concluded from this that proof
assistants provide the dense feedback that reinforcement learning requires. This
work reports what we found when we built such an environment and measured it. A
kernel makes the reward \emph{honest}, but it does nothing to make the reward
\emph{informative}, and these are different properties. We construct an
environment for information-flow security proofs in Lean~4, in which difficulty
is controlled by how much of a lemma's dependency structure is withheld from the
model, and we observe that two of our three configurations produce a constant
reward and therefore no learning signal at all, while appearing entirely healthy
from the outside. We further show that the difficulty can be moved deliberately,
and that the mechanism for moving it is decomposition: splitting a proof into
intermediate lemmas took a theorem from unprovable to provable for a fixed model
at a fixed budget. Along the way we report that the dependency structure on
which the whole construction rests cannot be read from the Lean environment at
all, and give the alternative we used instead.
\end{abstract}

\tableofcontents

\section{Introduction}

\subsection{Motivation}

A recurring proposal in machine learning for programs is to train models with
reinforcement learning against a reward derived from a proof assistant. The
appeal is that the reward is ground truth. When a model is trained against a
test suite, passing tests is a proxy for correctness, and it is a poor one: test
coverage is not strongly related to the number of defects a program contains, so
a model may improve its reward without improving the property one cares about.
A proof assistant does not have this weakness. Its kernel decides whether a
candidate proof establishes the stated theorem, and no amount of optimisation
pressure changes that decision.

From this it is frequently concluded that verification supplies the dense
feedback signal that reinforcement learning needs. The purpose of this work is
to examine that conclusion by building the environment and measuring it.

\subsection{Background}

It is useful to separate three things that are often discussed together.

A \emph{proof corpus} is a collection of statements and proofs. It is not an
environment: it contains no notion of difficulty, it cannot generate tasks that
have not been written, and if it is public then a model may have seen it during
training.

A \emph{kernel} is a decision procedure for proof validity. It is not an
environment either: it scores a candidate that is presented to it, but it says
nothing about which candidates should be presented.

An \emph{environment}, in the sense required for reinforcement learning, is a
distribution over tasks together with a reward function, such that a policy
interacting with it receives a signal that varies with its behaviour. The last
clause is the one that matters here, and it is the one that a kernel does not
supply.

\subsection{Objective}

We set out to build an environment of the third kind for a specific domain,
information-flow security, and then to characterise it: to determine whether the
tasks it produces actually elicit a varying reward from contemporary models, and
if not, to identify what governs that.

We chose information-flow security because it admits a family of related
theorems of graded difficulty. A security type system assigns confidentiality
levels to variables and constrains which assignments are permitted, and the
theorem that such a system is sound, called non-interference, is proved through a
chain of supporting lemmas. That chain gives the environment its structure.

\subsection{Contributions}

\begin{itemize}
\item We construct an environment for Lean~4 proof tasks in which difficulty is
controlled by a single parameter, namely how much of a lemma's dependency cone
is withheld, and in which the reward is kernel acceptance together with an audit
of the axioms the result depends on.
\item We report that the dependency structure required to define that parameter
cannot be obtained from the Lean environment, because proof terms are discarded
after checking, and we give a replacement procedure based on deletion.
\item We measure three configurations of the environment and find that two
produce a constant reward, and therefore no learning signal, while presenting no
outward sign of doing so.
\item We show that difficulty is controllable, and that the control is the
granularity of decomposition: a theorem that a model could not prove became
provable, at the same budget, once its proof was split into intermediate lemmas.
\item We observe that the environment admits a free estimate of its own
measurement noise, and report that this noise is large enough to account for
several apparent effects.
\end{itemize}

\section{The environment}

\subsection{Ladders and dependency depth}

The basic object is what we call a \emph{ladder}: a Lean theory together with a
set of lemmas proved over it, and the dependency relation between those lemmas.

\begin{definition}
The \emph{depth} of a lemma $L$ is $1$ if $L$ has no dependencies, and otherwise
$1 + \max\{\mathrm{depth}(A) : A \text{ is a dependency of } L\}$.
\end{definition}

Depth is computed from the dependency relation rather than assigned by hand.
This is deliberate, and Section~\ref{sec:extract} explains why it must be.

For example, in the security type system we use, the fact that agreement of two
memories on public variables is a transitive relation has depth~$1$, since it
follows directly from transitivity of equality. The confinement lemma, which
states that a command typed in a secret context cannot modify the public part of
the memory, has depth~$2$, because its proof requires the transitivity fact
together with two others. The soundness theorem itself has depth~$5$.

\subsection{Tasks and the difficulty parameter}

A \emph{task} presents the model with a Lean file and asks it to complete one or
more proofs. What varies between tasks is how much of the surrounding
development it is given.

\begin{definition}
Given a target lemma $L$, a task is determined by a set $W$ of \emph{withheld}
lemmas, where $L \in W$ and $W$ is contained in the dependency cone of $L$.
Lemmas in $W$ appear as goals the model must prove. Lemmas in the cone but not
in $W$ appear in the file already proved, and may be cited.
\end{definition}

Two extremes are of particular interest. When $W = \{L\}$, every dependency is
supplied to the model as an already-established result; we call this the
\emph{supplied} condition. When $W$ is the entire cone, the model must
reconstruct the whole subtree beneath $L$ before it can prove $L$ itself; we call
this the \emph{withheld} condition. At depth $5$ this is the difference between
writing three lines and writing eight lemmas.

Intermediate choices of $W$ interpolate between these, which is what makes
difficulty a parameter rather than a property of the corpus. It is also worth
observing that a supplied lemma is exactly what is meant by a trusted interface
in compositional verification, so this parameter is not an artificial construct
but the thing that architecture is about.

\subsection{Reward}

A candidate is accepted if three conditions hold. The file must compile. The
declaration \verb|#print axioms| applied to each target must report a subset of
the three axioms that Lean's own standard library depends upon. And the model's
text must not contain certain tactics, principally \verb|native_decide|, which
delegates to the compiler rather than the kernel.

The second condition is not redundant with the first, and this is worth stating
because we did not anticipate it. During development, one lemma failed to prove.
Lean reported the error, and nonetheless \emph{declared the theorem}. Every
later lemma that cited it compiled with exit status zero and produced no
diagnostic of its own. The only remaining trace, three lemmas further on, was:

\begin{lstlisting}
'L07' depends on axioms: [propext, sorryAx]
\end{lstlisting}

A reward function keyed on successful compilation would have paid out for that
entire subtree. We therefore treat the axiom audit, and not the exit status, as
the acceptance criterion.

\subsection{Preventing the model from changing the question}

A model asked to produce a proof may instead produce a proof of something
weaker. We remove this possibility structurally rather than by detection. The
harness emits the theorem statement itself, in the form

\begin{lstlisting}
theorem <name> <binders> : <statement> := by
\end{lstlisting}

and the model supplies only the tactic block that follows. Since the model never
writes the statement, it cannot weaken it. We had first implemented a comparison
between the model's statement and the intended one, and removed it: making a
failure impossible is preferable to detecting it.

\section{Recovering the dependency structure}
\label{sec:extract}

The difficulty parameter of Section~2.2 is defined in terms of a lemma's
dependency cone, so the environment cannot be built without knowing, for each
lemma, which other lemmas its proof requires. This turned out to be the hardest
part of the construction.

\subsection{The direct approach is unavailable}

The natural method is to ask Lean. The kernel has just checked the proof, so it
has seen every constant the proof refers to. In Lean~4 this information is not
retained. Theorem bodies are elaborated asynchronously and the resulting term is
not kept, so the field that should contain it is empty:

\begin{lstlisting}
theorem foo (n : Nat) : n = n := rfl
theorem bar (n : Nat) : n = n := foo n     -- cites foo

env.find? `bar                  -- value?.isSome = false
env.toKernelEnv.find? `bar      -- value?.isSome = false
liftCoreM (getConstInfo `bar)   -- value?.isSome = false
liftCoreM (collectAxioms `bar)  -- []
\end{lstlisting}

Two lines after \verb|bar| is declared, its proof is already unreachable.
Compiling to a \verb|.olean| file and loading it through
\verb|withImportModules| gives the same result, so this is not an artefact of an
unflushed in-memory structure. Only the constants occurring in the
\emph{statement} survive. Applied to our largest development, extraction by this
route reported all $46$ theorems as having no dependencies, and therefore a
maximum depth of $1$.

\subsection{Dependency by necessity}

We therefore determine the relation semantically. We assert that $L$ depends on
$A$ precisely when removing $A$ from the file causes $L$'s proof to fail.

This is not merely a workaround; it is a stronger criterion than the one we could
not use. A proof may mention a lemma without that lemma being load-bearing, in
which case reading the proof term would record a dependency that does not exist.
What the difficulty parameter requires is necessity, and necessity is what the
deletion test measures.

One complication must be handled. Removing $A$ also breaks any supplied lemma
whose own proof requires $A$, and those failures would otherwise be attributed to
$L$. We therefore process lemmas in order, so that the relation among earlier
lemmas is already known, and remove the whole set of lemmas that depend on $A$
together with $A$ itself. A failure then establishes that $L$ requires something
in that set, hence that $L$ transitively requires $A$. Direct dependencies are
recovered afterwards by transitive reduction.

\subsection{Validation}

We had, at the time, a dependency relation that had been written by hand.
Applying the procedure to it produced the following:

\begin{center}
\begin{tabular}{lll}
\toprule
asserted dependency & established by deletion? & verdict \\
\midrule
\verb|lowEq_symm| requires \verb|lowEq_refl|   & no & rejected \\
\verb|lowEq_trans| requires \verb|lowEq_refl|  & no & rejected \\
\verb|Lvl.le_trans| requires \verb|Lvl.le_refl|& no & rejected \\
\verb|secure_sub| requires \verb|Lvl.le_refl|  & no & rejected \\
\verb|sub_base_inv| requires \verb|Lvl.le_refl|& yes & accepted \\
\verb|confinement| requires \verb|secure_sub|  & yes & accepted \\
\bottomrule
\end{tabular}
\end{center}

Four of the fourteen asserted dependencies did not exist. The first of them
claimed that symmetry of memory agreement requires reflexivity of the same
relation, whereas its proof is

\begin{lstlisting}
theorem lowEq_symm {G : Ctx} {s t : St} (h : lowEq G s t) : lowEq G t s := by
  intro x hx; exact (h x hx).symm
\end{lstlisting}

which cannot use reflexivity, since reflexivity relates a memory to itself and
says nothing about exchanging two memories. Deleting the reflexivity lemma from
the file leaves this proof compiling and axiom-clean.

The consequence was that the reported maximum depth of that ladder was $6$ when
the true value is $5$, and every measurement indexed by depth had been taken
against the wrong scale. We note that a check intended to detect exactly this had
been reporting no failures throughout, and was on inspection incapable of
reporting anything else: it inferred that a dependency was genuine from a
rejection that the harness itself had caused. We now require that every check in
the system be demonstrated failing on a case it is meant to reject before it is
relied upon.

\section{The corpus}

The environment contains fourteen theories: twelve generated over a parameter
space, one written by hand as a control, and one that is a full formalisation.
Together they contain $158$ lemmas and yield $2{,}607$ distinct tasks. Every
dependency edge was established by the procedure of Section~3, none by hand.

The formalisation is a security type system in the style of Volpano, Smith and
Irvine, extended with subtyping, while loops, and functions. It is $784$ lines of
Lean, contains no \verb|sorry|, and its soundness theorem depends only on
\verb|propext| and \verb|Quot.sound|. Because evaluation of a language with loops
need not terminate, the semantics is indexed by a fuel parameter, and the theorem
proved is accordingly termination-insensitive: it relates two executions that
both complete. A structural consequence is that agreement of expressions and
non-interference of commands must be proved simultaneously, by induction on the
fuel, since a function call occurring in an expression executes a command while
an assignment occurring in a command evaluates an expression.

\section{Measurements}

All measurements use a single attempt per task, with no access to a compiler, so
the model cannot iterate against the kernel. As controls, the reference solutions
score $1.00$ in every cell and a trivial tactic scores $0.00$, which confirms
that the reward discriminates.

\subsection{The kernel does not supply a learning signal}

We ran three configurations of ladder and model. The following gives the
proportion of tasks solved at each depth in the withheld condition.

\begin{center}
\begin{tabular}{llc}
\toprule
ladder & model & pass rate at depths $1,2,3,4$ \\
\midrule
arithmetic & Sonnet~5 & $1.00,\ 1.00,\ 1.00,\ 1.00$ \\
security types, first version & Sonnet~5 & $1.00,\ 0.00,\ 0.00,\ 0.00$ \\
arithmetic & Haiku~4.5 & $1.00,\ 0.75,\ 0.33,\ 0.00$ \\
\bottomrule
\end{tabular}
\end{center}

In the first configuration the model solves everything, and in the second it
solves nothing beyond the trivial lemmas. In both, the reward is constant. A
policy interacting with either receives no information about which of its
behaviours were better, which is precisely the sparse-reward condition that the
appeal to proof assistants was supposed to avoid. Only the third configuration
produces a reward that varies with depth.

We stress that nothing distinguishes the first two from the outside. Every task
is well formed, every reward is correctly computed, the kernel is behaving
exactly as intended, and the test suite reports no failures. An environment that
teaches nothing looks identical to one that teaches, unless one measures the
distribution of rewards it actually produces.

We therefore state the conclusion as follows. Dense feedback is not a property
conferred by a proof assistant. It is a property of a task distribution
considered together with a particular policy. The kernel guarantees that the
reward is honest; it makes no guarantee that the reward is informative.

\subsection{Difficulty is controllable, and decomposition is the control}

The second configuration above is not merely a bad choice of theory; it was
repaired, and the repair is instructive.

That ladder failed above depth~$1$ because its soundness theorem was proved in a
single step of thirty-seven lines, with nothing between it and the elementary
lemmas. We divided that proof into two intermediate lemmas, one handling the
assignment case and one handling the conditional with a secret guard, which
reduced the final proof to twenty-two lines. The same model, given the same
single attempt and no other assistance, then proved the theorem it had
previously failed.

No mathematics changed. The theory, the statement, and the model were identical;
only the granularity at which intermediate results were made available differed.
We regard this as the claim of compositional verification appearing as a property
of the environment: whether a corpus is learnable at all depends on how finely it
is cut.

Measured over the ladder as a whole, supplying a lemma's dependencies rather than
withholding them raises the proportion solved from $4/14$ to $18/27$, which is
significant at $p = 0.026$ by Fisher's exact test. We report this for
completeness and would not rest an argument on it; the sample is small and the
result is sensitive to how repeated responses are counted.

\subsection{The environment measures its own noise}

A lemma of depth~$1$ has no dependencies. Consequently the supplied and withheld
conditions, applied to such a lemma, produce files that are identical. We
verified this with \verb|diff|. They are nevertheless dispatched to independent
instances of the model under different labels.

On one such set, the two arms returned $8/8$ and $6/8$ on identical input. The
implied variation between runs is approximately twenty-five percent.

This is a useful accident. Any environment built in this way contains tasks that
are duplicates by construction, and therefore obtains an estimate of its own
measurement noise at no additional cost. We recommend including such tasks
deliberately. In our own case the estimate was necessary: without it we would
have read several single-sample cells as effects, and one of our earlier reported
results was of exactly that kind.

\section{Discussion}

The components we expected to be difficult were not. The reward function, the
task generator, the renaming that protects against memorisation, and the
machinery that establishes dependencies all worked once built, and none of them
determined whether the environment was useful.

What determined it was calibration, in the specific sense of whether the tasks
are hard enough that the model does not solve them all and easy enough that it
does not fail them all. This has no natural alarm. A ceiling and a floor look the
same as each other and the same as a healthy environment, from every vantage
point except the distribution of rewards.

We would put the practical recommendation as follows. If one intends to scale
reinforcement learning on program verification, the proof assistant is the part
of the problem that is already solved. What is required in addition is a task
distribution whose difficulty can be adjusted, instrumentation that detects when
the reward has become constant, and an estimate of measurement noise against
which apparent effects can be judged. We have found the third of these to be
obtainable for free, and the first to be obtainable through decomposition.

\section*{Availability}

The environment, the corpus, the extraction procedure, the formalisation, and a
record of the errors described above are available at
\url{https://github.com/ymiled/proof-depth-decay}.

\end{document}
